\documentclass[11pt,a4paper]{article}

% ---- Packages ----
\usepackage[utf8]{inputenc}
\usepackage[T1]{fontenc}
\usepackage{amsmath,amssymb,amsthm}
\usepackage{mathtools}
\usepackage[margin=1in,headheight=14pt]{geometry}
\usepackage{hyperref}
\usepackage{enumitem}
\usepackage{xcolor}
\usepackage{tcolorbox}
\usepackage{fancyhdr}
\usepackage{booktabs}

% ---- Hyperref ----
\hypersetup{
    colorlinks=true,
    linkcolor=red,
    citecolor=blue,
    urlcolor=blue
}

% ---- Header ----
\pagestyle{fancy}
\fancyhf{}
\fancyhead[L]{\textit{Probability Foundations for Reinforcement Learning}}
\fancyhead[R]{\thepage}
\renewcommand{\headrulewidth}{0.4pt}

% ---- Theorem Environments ----
\theoremstyle{definition}
\newtheorem{definition}{Definition}[section]
\newtheorem{assumption}{Assumption}[section]
\newtheorem{example}{Example}[section]

\theoremstyle{plain}
\newtheorem{theorem}{Theorem}[section]
\newtheorem{lemma}[theorem]{Lemma}
\newtheorem{proposition}[theorem]{Proposition}
\newtheorem{corollary}[theorem]{Corollary}

\theoremstyle{remark}
\newtheorem{remark}{Remark}[section]

% ---- Colored Boxes ----
\tcbuselibrary{theorems,skins,breakable}

\newtcolorbox{keyresult}[1][]{
    colback=green!5!white,
    colframe=green!50!black,
    fonttitle=\bfseries,
    title=#1,
    breakable
}

\newtcolorbox{rlconnection}[1][]{
    colback=blue!5!white,
    colframe=blue!50!black,
    fonttitle=\bfseries,
    title=#1,
    breakable
}

\newtcolorbox{tdconnection}[1][]{
    colback=blue!5!white,
    colframe=blue!50!black,
    fonttitle=\bfseries,
    title=#1,
    breakable
}

% ---- Operators ----
\DeclareMathOperator{\Var}{Var}
\DeclareMathOperator{\Cov}{Cov}
\newcommand{\E}{\mathbb{E}}
\newcommand{\R}{\mathbb{R}}
\newcommand{\N}{\mathbb{N}}
\newcommand{\Prob}{\mathbb{P}}
\newcommand{\indicator}{\mathbf{1}}
\newcommand{\cas}{\xrightarrow{\text{a.s.}}}
\newcommand{\cprob}{\xrightarrow{\Prob}}
\newcommand{\clp}[1]{\xrightarrow{L^{#1}}}
\newcommand{\cdist}{\xrightarrow{d}}

% ---- Title ----
\title{Lesson 5: Probability Spaces, Random Variables, and Expectation}
\author{AI/ML Mathematics Series}
\date{February 2026}

\begin{document}
\maketitle
\tableofcontents
\newpage

\setcounter{section}{0}

\section{Introduction}
\label{sec:intro}

The convergence proofs in reinforcement learning rely on a core set of probabilistic
tools. This document develops the complete probability foundation from first
principles, providing every definition, theorem, and proof needed to follow the
convergence analyses in the companion articles on Temporal-Difference learning,
policy gradient methods, and stochastic approximation.

\medskip\noindent\textbf{What this document covers.}
\begin{enumerate}[nosep]
    \item \textbf{Probability spaces and events} --- $\sigma$-algebras, Kolmogorov axioms, conditional probability, Bayes' theorem.
    \item \textbf{Random variables and distributions} --- measurability, distribution functions, expectation, common distributions, joint distributions, linearity.
    \item \textbf{Conditional expectation and the tower property} --- conditional expectation given events and random variables, linearity, pull-out property, law of total expectation.
    \item \textbf{MDP expectation expansion} --- the mechanical expansion step in Bellman derivations.
    \item \textbf{Variance and conditional variance} --- computational formulas, law of total variance.
    \item \textbf{Independence} --- independent events, independent random variables, product of expectations.
    \item \textbf{Filtrations} --- natural filtration of an MDP, the Markov property and conditioning equivalence.
    \item \textbf{Martingale difference sequences} --- definitions, zero mean, partial sums form martingales, uncorrelated increments.
    \item \textbf{Probability inequalities} --- Markov, Chebyshev, Cauchy--Schwarz, Hoeffding, Azuma--Hoeffding, Jensen.
    \item \textbf{Modes of convergence} --- almost sure, in probability, in $L^p$, in distribution, and their relationships.
    \item \textbf{Borel--Cantelli lemmas} --- the bridge between summable probabilities and almost sure convergence.
    \item \textbf{Limit theorems} --- weak and strong law of large numbers, central limit theorem, monotone and dominated convergence.
    \item \textbf{Markov chains} --- transition matrices, stationary distributions, ergodicity.
    \item \textbf{Martingale convergence theory} --- supermartingales, Doob's convergence theorem, Robbins--Siegmund.
    \item \textbf{Stochastic approximation noise decomposition} --- the TD noise pattern that drives all convergence proofs.
\end{enumerate}

%% ============================================================
%% SECTION 2: PROBABILITY SPACES AND EVENTS
%% ============================================================

\section{Probability Spaces and Events}
\label{sec:prob_spaces}

\subsection{Sample Spaces and \texorpdfstring{$\sigma$}{Sigma}-Algebras}
\label{subsec:sigma_algebras}

\begin{definition}[$\sigma$-Algebra]
\label{def:sigma_algebra}
A collection $\mathcal{F}$ of subsets of a set $\Omega$ is a \emph{$\sigma$-algebra}
(or \emph{$\sigma$-field}) if:
\begin{enumerate}[nosep]
    \item $\Omega \in \mathcal{F}$,
    \item If $A \in \mathcal{F}$, then $A^c = \Omega \setminus A \in \mathcal{F}$ \hfill(closed under complement),
    \item If $A_1, A_2, \ldots \in \mathcal{F}$, then $\bigcup_{n=1}^\infty A_n \in \mathcal{F}$ \hfill(closed under countable union).
\end{enumerate}
\end{definition}

\begin{proposition}[Basic Properties of $\sigma$-Algebras]
\label{prop:sigma_props}
Let $\mathcal{F}$ be a $\sigma$-algebra on $\Omega$. Then:
\begin{enumerate}[nosep]
    \item $\emptyset \in \mathcal{F}$.
    \item $\mathcal{F}$ is closed under countable intersection: if $A_1, A_2, \ldots \in \mathcal{F}$, then $\bigcap_{n=1}^\infty A_n \in \mathcal{F}$.
    \item $\mathcal{F}$ is closed under set difference: if $A, B \in \mathcal{F}$, then $A \setminus B \in \mathcal{F}$.
\end{enumerate}
\end{proposition}

\begin{proof}
\textbf{(1)} $\emptyset = \Omega^c \in \mathcal{F}$ by closure under complement.

\textbf{(2)} By De Morgan's law, $\bigcap_{n=1}^\infty A_n = \left(\bigcup_{n=1}^\infty A_n^c\right)^c$. Each $A_n^c \in \mathcal{F}$ by (2) of Definition~\ref{def:sigma_algebra}, so $\bigcup A_n^c \in \mathcal{F}$ by (3), and $\left(\bigcup A_n^c\right)^c \in \mathcal{F}$ by (2).

\textbf{(3)} $A \setminus B = A \cap B^c$. Since $B^c \in \mathcal{F}$, this is a finite intersection of elements of $\mathcal{F}$, which is a special case of (2).
\end{proof}

\begin{definition}[Generated $\sigma$-Algebra]
\label{def:generated_sigma}
Given a collection $\mathcal{C}$ of subsets of $\Omega$, the \emph{$\sigma$-algebra generated by $\mathcal{C}$}, denoted $\sigma(\mathcal{C})$, is the smallest $\sigma$-algebra containing $\mathcal{C}$:
\begin{equation}
    \sigma(\mathcal{C}) = \bigcap \{ \mathcal{G} : \mathcal{G} \text{ is a $\sigma$-algebra on } \Omega \text{ and } \mathcal{C} \subseteq \mathcal{G} \}.
    \label{eq:generated_sigma}
\end{equation}
\end{definition}

\begin{remark}
The intersection of any collection of $\sigma$-algebras is again a $\sigma$-algebra (verify the three axioms directly), so $\sigma(\mathcal{C})$ is well defined. The power set $2^\Omega$ is always a $\sigma$-algebra containing $\mathcal{C}$, so the intersection is non-empty.
\end{remark}

\begin{rlconnection}[Connection to RL]
In RL, $\sigma(S_0, A_0, R_1, \ldots, S_t)$ is the $\sigma$-algebra generated by the trajectory up to time $t$. It represents ``all information available at time $t$'' and is precisely the filtration $\mathcal{F}_t$ used throughout the TD convergence proofs.
\end{rlconnection}

\subsection{Probability Measures}
\label{subsec:prob_measures}

\begin{definition}[Probability Measure --- Kolmogorov Axioms]
\label{def:prob_measure}
A \emph{probability measure} on a measurable space $(\Omega, \mathcal{F})$ is a function $\Prob: \mathcal{F} \to [0,1]$ satisfying:
\begin{enumerate}[nosep]
    \item \textbf{Non-negativity:} $\Prob(A) \geq 0$ for all $A \in \mathcal{F}$.
    \item \textbf{Normalization:} $\Prob(\Omega) = 1$.
    \item \textbf{Countable additivity:} For pairwise disjoint events $A_1, A_2, \ldots \in \mathcal{F}$:
    \begin{equation}
        \Prob\!\left(\bigcup_{n=1}^\infty A_n\right) = \sum_{n=1}^\infty \Prob(A_n).
        \label{eq:countable_additivity}
    \end{equation}
\end{enumerate}
The triple $(\Omega, \mathcal{F}, \Prob)$ is called a \emph{probability space}.
\end{definition}

\begin{theorem}[Basic Properties of Probability]
\label{thm:prob_props}
Let $(\Omega, \mathcal{F}, \Prob)$ be a probability space. For events $A, B \in \mathcal{F}$:
\begin{enumerate}[nosep]
    \item $\Prob(\emptyset) = 0$.
    \item \textbf{Complement:} $\Prob(A^c) = 1 - \Prob(A)$.
    \item \textbf{Monotonicity:} If $A \subseteq B$, then $\Prob(A) \leq \Prob(B)$.
    \item \textbf{Inclusion-exclusion:} $\Prob(A \cup B) = \Prob(A) + \Prob(B) - \Prob(A \cap B)$.
    \item \textbf{Sub-additivity:} $\Prob(A \cup B) \leq \Prob(A) + \Prob(B)$.
    \item \textbf{Union bound:} $\Prob\!\left(\bigcup_{n=1}^\infty A_n\right) \leq \sum_{n=1}^\infty \Prob(A_n)$.
\end{enumerate}
\end{theorem}

\begin{proof}
\textbf{(1)} Write $\Omega = \Omega \cup \emptyset \cup \emptyset \cup \cdots$ (disjoint). By countable additivity, $\Prob(\Omega) = \Prob(\Omega) + \Prob(\emptyset) + \Prob(\emptyset) + \cdots$. Since $\Prob(\Omega) = 1 < \infty$, all remaining terms must be zero, so $\Prob(\emptyset) = 0$.

\textbf{(2)} $\Omega = A \cup A^c$ with $A \cap A^c = \emptyset$, so $1 = \Prob(A) + \Prob(A^c)$.

\textbf{(3)} Write $B = A \cup (B \setminus A)$ disjointly. Then $\Prob(B) = \Prob(A) + \Prob(B \setminus A) \geq \Prob(A)$.

\textbf{(4)} Write $A \cup B = A \cup (B \setminus A)$ disjointly, and $B = (A \cap B) \cup (B \setminus A)$ disjointly. From the second: $\Prob(B \setminus A) = \Prob(B) - \Prob(A \cap B)$. Substituting into the first: $\Prob(A \cup B) = \Prob(A) + \Prob(B) - \Prob(A \cap B)$.

\textbf{(5)} From (4), since $\Prob(A \cap B) \geq 0$.

\textbf{(6)} Define $B_1 = A_1$ and $B_n = A_n \setminus \bigcup_{k=1}^{n-1} A_k$ for $n \geq 2$. The $B_n$ are pairwise disjoint with $\bigcup B_n = \bigcup A_n$, and $B_n \subseteq A_n$ so $\Prob(B_n) \leq \Prob(A_n)$. Then $\Prob(\bigcup A_n) = \Prob(\bigcup B_n) = \sum \Prob(B_n) \leq \sum \Prob(A_n)$.
\end{proof}

\begin{theorem}[Continuity of Probability]
\label{thm:continuity}
Let $(\Omega, \mathcal{F}, \Prob)$ be a probability space.
\begin{enumerate}[nosep]
    \item \textbf{Continuity from below:} If $A_1 \subseteq A_2 \subseteq \cdots$, then $\Prob\!\left(\bigcup_{n=1}^\infty A_n\right) = \lim_{n \to \infty} \Prob(A_n)$.
    \item \textbf{Continuity from above:} If $A_1 \supseteq A_2 \supseteq \cdots$, then $\Prob\!\left(\bigcap_{n=1}^\infty A_n\right) = \lim_{n \to \infty} \Prob(A_n)$.
\end{enumerate}
\end{theorem}

\begin{proof}
\textbf{(1)} Define $B_1 = A_1$ and $B_n = A_n \setminus A_{n-1}$ for $n \geq 2$. We establish three claims.

\medskip
\textit{Claim 1: The $B_n$ are pairwise disjoint.}
Let $i < j$. We must show $B_i \cap B_j = \emptyset$, i.e., there is no $\omega$ belonging to both $B_i$ and $B_j$. We proceed by contradiction.

Suppose $\omega \in B_i \cap B_j$.

\emph{Step 1: $\omega \in B_i$ implies $\omega \in A_{j-1}$.}
If $i = 1$, then $B_1 = A_1$, so $\omega \in A_1$. If $i \geq 2$, then $B_i = A_i \setminus A_{i-1} \subseteq A_i$, so $\omega \in A_i$.
In either case $\omega \in A_i$. Since $i < j$ and the sets are nested ($A_1 \subseteq A_2 \subseteq \cdots$), we have $i \leq j - 1$, hence $A_i \subseteq A_{j-1}$. Therefore $\omega \in A_{j-1}$.

\emph{Step 2: $\omega \in B_j$ implies $\omega \notin A_{j-1}$.}
Since $j \geq 2$ (because $i \geq 1$ and $i < j$), we have $B_j = A_j \setminus A_{j-1} = \{\alpha \in A_j : \alpha \notin A_{j-1}\}$.
So $\omega \in B_j$ gives $\omega \notin A_{j-1}$ by definition of set difference.

\emph{Conclusion.}
Step~1 gives $\omega \in A_{j-1}$ and Step~2 gives $\omega \notin A_{j-1}$. This is a contradiction, so no such $\omega$ exists and $B_i \cap B_j = \emptyset$.

\medskip
\textit{Claim 2: $A_N = \bigcup_{k=1}^N B_k$ for every $N \geq 1$.}
We prove this by induction on $N$.

\emph{Base case} ($N = 1$): $\bigcup_{k=1}^1 B_k = B_1 = A_1$. \checkmark

\emph{Inductive step}: Assume $A_{N-1} = \bigcup_{k=1}^{N-1} B_k$. Then:
\begin{equation}
    \bigcup_{k=1}^N B_k = \underbrace{\bigcup_{k=1}^{N-1} B_k}_{= A_{N-1}} \cup\; B_N
    = A_{N-1} \cup (A_N \setminus A_{N-1}).
    \label{eq:disjointification_induction}
\end{equation}
For any sets $C \subseteq D$, we have $C \cup (D \setminus C) = D$ (every element of $D$ is either in $C$ or in $D \setminus C$). Since $A_{N-1} \subseteq A_N$, equation~\eqref{eq:disjointification_induction} gives $\bigcup_{k=1}^N B_k = A_N$. \checkmark

\medskip
\textit{Claim 3: $\sum_{n=1}^N \Prob(B_n) = \Prob(A_N)$.}
By Claim~1, $B_1, \ldots, B_N$ are pairwise disjoint. By finite additivity (which follows from countable additivity applied with $A_{N+1} = A_{N+2} = \cdots = \emptyset$):
\begin{equation}
    \sum_{n=1}^N \Prob(B_n) = \Prob\!\left(\bigcup_{k=1}^N B_k\right) \overset{\text{Claim 2}}{=} \Prob(A_N).
    \label{eq:finite_additivity_step}
\end{equation}

\medskip
Now, since $\bigcup_{n=1}^\infty A_n = \bigcup_{n=1}^\infty B_n$ (Claim~2 gives $\subseteq$; $B_n \subseteq A_n$ gives $\supseteq$), countable additivity applies to the disjoint union (Claim~1):
\begin{equation}
    \Prob\!\left(\bigcup_{n=1}^\infty A_n\right) = \Prob\!\left(\bigcup_{n=1}^\infty B_n\right) = \sum_{n=1}^\infty \Prob(B_n) = \lim_{N \to \infty} \sum_{n=1}^N \Prob(B_n) \overset{\eqref{eq:finite_additivity_step}}{=} \lim_{N \to \infty} \Prob(A_N).
\end{equation}

\textbf{(2)} Apply (1) to $A_n^c$: since $A_1 \supseteq A_2 \supseteq \cdots$, the complements satisfy $A_1^c \subseteq A_2^c \subseteq \cdots$, so by (1): $\Prob(\bigcup_{n=1}^\infty A_n^c) = \lim_{n \to \infty} \Prob(A_n^c)$. By De Morgan, $\bigcup A_n^c = (\bigcap A_n)^c$, giving $1 - \Prob(\bigcap A_n) = \lim(1 - \Prob(A_n))$, hence $\Prob(\bigcap A_n) = \lim \Prob(A_n)$.
\end{proof}

\begin{rlconnection}[Connection to RL]
Continuity from above is the key tool for defining ``almost sure'' events. The event ``$V_t$ converges to $V^\pi$'' is an intersection of increasingly restrictive events, and continuity from above lets us compute its probability as a limit. This underpins every ``convergence with probability 1'' result in the TD article.
\end{rlconnection}

\subsection{Conditional Probability and Bayes' Theorem}
\label{subsec:cond_prob}

\begin{definition}[Conditional Probability]
\label{def:cond_prob}
For events $A, B \in \mathcal{F}$ with $\Prob(B) > 0$, the \emph{conditional probability of $A$ given $B$} is:
\begin{equation}
    \Prob(A \mid B) = \frac{\Prob(A \cap B)}{\Prob(B)}.
    \label{eq:cond_prob}
\end{equation}
\end{definition}

\begin{proposition}[Conditional Probability is a Probability Measure]
\label{prop:cond_is_prob}
For a fixed event $B$ with $\Prob(B) > 0$, the function $A \mapsto \Prob(A \mid B)$ satisfies the Kolmogorov axioms and is therefore a probability measure on $(\Omega, \mathcal{F})$.
\end{proposition}

\begin{proof}
\textbf{Non-negativity:} $\Prob(A \mid B) = \Prob(A \cap B)/\Prob(B) \geq 0$.

\textbf{Normalization:} $\Prob(\Omega \mid B) = \Prob(\Omega \cap B)/\Prob(B) = \Prob(B)/\Prob(B) = 1$.

\textbf{Countable additivity:} For pairwise disjoint $A_1, A_2, \ldots$:
\begin{align}
    \Prob\!\left(\bigcup_n A_n \;\middle|\; B\right) &= \frac{\Prob\!\left(\left(\bigcup_n A_n\right) \cap B\right)}{\Prob(B)} = \frac{\Prob\!\left(\bigcup_n (A_n \cap B)\right)}{\Prob(B)} = \frac{\sum_n \Prob(A_n \cap B)}{\Prob(B)} = \sum_n \Prob(A_n \mid B).
\end{align}
\end{proof}

\begin{theorem}[Chain Rule of Probability]
\label{thm:chain_rule}
For events $A_1, \ldots, A_n$ with $\Prob(A_1 \cap \cdots \cap A_{n-1}) > 0$:
\begin{equation}
    \Prob(A_1 \cap A_2 \cap \cdots \cap A_n) = \Prob(A_1)\, \Prob(A_2 \mid A_1)\, \Prob(A_3 \mid A_1 \cap A_2) \cdots \Prob(A_n \mid A_1 \cap \cdots \cap A_{n-1}).
    \label{eq:chain_rule}
\end{equation}
\end{theorem}

\begin{proof}
By induction on $n$. The base case $n = 2$ is the definition: $\Prob(A_1 \cap A_2) = \Prob(A_1)\, \Prob(A_2 \mid A_1)$.

For the inductive step, assume the formula holds for $n-1$. Write $B = A_1 \cap \cdots \cap A_{n-1}$. Then $\Prob(B \cap A_n) = \Prob(B)\, \Prob(A_n \mid B)$, and expand $\Prob(B)$ by the inductive hypothesis.
\end{proof}

\begin{theorem}[Law of Total Probability]
\label{thm:total_prob}
Let $B_1, B_2, \ldots, B_n$ be a partition of $\Omega$ (i.e., pairwise disjoint with $\bigcup_i B_i = \Omega$) with $\Prob(B_i) > 0$ for all $i$. Then for any event $A$:
\begin{equation}
    \Prob(A) = \sum_{i=1}^n \Prob(A \mid B_i)\, \Prob(B_i).
    \label{eq:total_prob}
\end{equation}
\end{theorem}

\begin{proof}
Since the $B_i$ partition $\Omega$, we have $A = \bigcup_{i=1}^n (A \cap B_i)$ with the terms pairwise disjoint. By finite additivity and the definition of conditional probability:
\begin{equation}
    \Prob(A) = \sum_{i=1}^n \Prob(A \cap B_i) = \sum_{i=1}^n \Prob(A \mid B_i)\, \Prob(B_i).
\end{equation}
\end{proof}

\begin{keyresult}[Bayes' Theorem]
\begin{theorem}[Bayes' Theorem]
\label{thm:bayes}
Let $B_1, \ldots, B_n$ be a partition of $\Omega$ with $\Prob(B_i) > 0$, and let $A$ be an event with $\Prob(A) > 0$. Then:
\begin{equation}
    \Prob(B_j \mid A) = \frac{\Prob(A \mid B_j)\, \Prob(B_j)}{\sum_{i=1}^n \Prob(A \mid B_i)\, \Prob(B_i)}.
    \label{eq:bayes}
\end{equation}
\end{theorem}
\end{keyresult}

\begin{proof}
By the definition of conditional probability and the law of total probability:
\begin{equation}
    \Prob(B_j \mid A) = \frac{\Prob(A \cap B_j)}{\Prob(A)} = \frac{\Prob(A \mid B_j)\, \Prob(B_j)}{\sum_{i=1}^n \Prob(A \mid B_i)\, \Prob(B_i)}.
\end{equation}
\end{proof}

\begin{rlconnection}[Connection to RL]
The chain rule of probability decomposes the joint probability of a trajectory $(S_0, A_0, R_1, S_1, \ldots)$ into a product of transition and policy probabilities. This decomposition is the starting point of the importance sampling ratio used in off-policy methods (TRPO, PPO). Bayes' theorem underlies posterior updates in Bayesian RL and model-based methods.
\end{rlconnection}

%% ============================================================
%% SECTION 3: RANDOM VARIABLES AND DISTRIBUTIONS
%% ============================================================

\section{Random Variables and Distributions}
\label{sec:random_variables}

\subsection{Measurability}
\label{subsec:measurability}

\begin{definition}[Measurable Function / Random Variable]
\label{def:measurable}
Let $(\Omega, \mathcal{F})$ be a measurable space. A function $X: \Omega \to \R$ is \emph{$\mathcal{F}$-measurable} (or a \emph{random variable}) if for every Borel set $B \subseteq \R$:
\begin{equation}
    X^{-1}(B) = \{\omega \in \Omega : X(\omega) \in B\} \in \mathcal{F}.
    \label{eq:measurable}
\end{equation}
Equivalently, it suffices to check that $\{X \leq x\} \in \mathcal{F}$ for all $x \in \R$.
\end{definition}

\begin{proposition}[Algebra of Measurable Functions]
\label{prop:measurable_algebra}
If $X$ and $Y$ are $\mathcal{F}$-measurable and $a, b \in \R$, then the following are $\mathcal{F}$-measurable:
\begin{enumerate}[nosep]
    \item $aX + bY$ (linear combinations),
    \item $XY$ (products),
    \item $|X|$ (absolute value),
    \item $\max(X, Y)$ and $\min(X, Y)$,
    \item $g(X)$ for any Borel-measurable function $g: \R \to \R$.
\end{enumerate}
\end{proposition}

\begin{proof}
We prove (1); the others follow similar arguments.

For $aX + bY$, we need $\{aX + bY \leq c\} \in \mathcal{F}$ for all $c \in \R$. Note that $\{aX + bY \leq c\} = \bigcup_{q \in \mathbb{Q}} (\{aX \leq q\} \cap \{bY \leq c - q\})$. For $a > 0$, $\{aX \leq q\} = \{X \leq q/a\} \in \mathcal{F}$; similarly for $bY$. The countable union over $q \in \mathbb{Q}$ of sets in $\mathcal{F}$ is in $\mathcal{F}$.
\end{proof}

\begin{definition}[$\sigma$-Algebra Generated by a Random Variable]
\label{def:sigma_rv}
The $\sigma$-algebra generated by a random variable $X$ is:
\begin{equation}
    \sigma(X) = \{X^{-1}(B) : B \subseteq \R \text{ is Borel}\} = \{\{X \in B\} : B \text{ Borel}\}.
    \label{eq:sigma_rv}
\end{equation}
A random variable $Y$ is $\sigma(X)$-measurable if and only if $Y = g(X)$ for some Borel function $g$.
\end{definition}

\begin{remark}
The statement ``$Y$ is $\sigma(X)$-measurable if and only if $Y = g(X)$'' is called the \emph{Doob--Dynkin lemma}. It formalizes the intuitive idea that $Y$ is ``determined by $X$'' exactly when $Y$ is a function of $X$. This is the foundation of the pull-out property (Theorem~\ref{thm:pull_out} below): if $h(Z)$ is $\sigma(Z)$-measurable, it acts as a constant when we condition on $Z$.
\end{remark}

\subsection{Distribution Functions}
\label{subsec:distributions}

\begin{definition}[Cumulative Distribution Function]
\label{def:cdf}
The \emph{cumulative distribution function} (CDF) of a random variable $X$ is:
\begin{equation}
    F_X(x) = \Prob(X \leq x), \quad x \in \R.
    \label{eq:cdf}
\end{equation}
\end{definition}

\begin{proposition}[Properties of the CDF]
\label{prop:cdf_props}
Every CDF $F$ satisfies:
\begin{enumerate}[nosep]
    \item $F$ is non-decreasing.
    \item $\lim_{x \to -\infty} F(x) = 0$ and $\lim_{x \to +\infty} F(x) = 1$.
    \item $F$ is right-continuous: $\lim_{y \downarrow x} F(y) = F(x)$.
\end{enumerate}
\end{proposition}

\begin{proof}
\textbf{(1)} If $x \leq y$, then $\{X \leq x\} \subseteq \{X \leq y\}$, so $F(x) = \Prob(X \leq x) \leq \Prob(X \leq y) = F(y)$ by monotonicity of $\Prob$.

\textbf{(2)} For the left limit, let $A_n = \{X \leq -n\} = \{\omega \in \Omega : X(\omega) \leq -n\}$.

\emph{The $A_n$ are decreasing:} If $\omega \in A_{n+1}$, then $X(\omega) \leq -(n+1) < -n$, so $\omega \in A_n$. Hence $A_{n+1} \subseteq A_n$, giving $A_1 \supseteq A_2 \supseteq \cdots$.

\emph{$\bigcap_{n=1}^\infty A_n = \emptyset$:} Suppose for contradiction that $\omega \in \bigcap_{n=1}^\infty A_n$. Then $\omega \in A_n$ for every $n \geq 1$, meaning $X(\omega) \leq -n$ for every $n \geq 1$. Since $X : \Omega \to \mathbb{R}$, the value $X(\omega)$ is a real number. But for any $x \in \mathbb{R}$, the Archimedean property of $\mathbb{R}$ guarantees the existence of $n_0 \in \mathbb{N}$ with $n_0 > -x$, i.e., $x > -n_0$. Taking $x = X(\omega)$, we get $X(\omega) > -n_0$, contradicting $X(\omega) \leq -n_0$. Therefore no such $\omega$ exists and $\bigcap_{n=1}^\infty A_n = \emptyset$.

By continuity from above (Theorem~\ref{thm:continuity}), $\lim_{n \to \infty} F(-n) = \lim_{n \to \infty} \Prob(A_n) = \Prob(\bigcap A_n) = \Prob(\emptyset) = 0$.

The right limit is analogous with $A_n = \{X \leq n\}$ increasing to $\Omega$ (every $\omega$ satisfies $X(\omega) \leq n$ for large enough $n$, again by the Archimedean property).

\textbf{(3)} Let $y_n \downarrow x$. Then $\{X \leq y_n\} \supseteq \{X \leq y_{n+1}\}$ and $\bigcap_n \{X \leq y_n\} = \{X \leq x\}$. By continuity from above, $\lim F(y_n) = F(x)$.
\end{proof}

\subsection{Expectation}
\label{subsec:expectation}

\begin{definition}[Expectation]
\label{def:expectation}
The \emph{expectation} (or \emph{expected value}) of a discrete random variable $X$ with PMF $p_X$ is
\begin{equation}
    \E[X] = \sum_{x} x\, p_X(x),
    \label{eq:expectation_def}
\end{equation}
provided the sum converges absolutely: $\sum_x |x|\, p_X(x) < \infty$.
\end{definition}

\begin{tdconnection}[Connection to TD Learning]
The state-value function $V^\pi(s) = \E_\pi[G_0 \mid S_0 = s]$ is an expectation of the return $G_0 = \sum_{t=0}^\infty \gamma^t R_{t+1}$ over all trajectories from state $s$. Since the MDP has finite state and action spaces and $\gamma < 1$, the return is bounded, so the expectation is always well defined.
\end{tdconnection}

\subsection{Common Distributions}
\label{subsec:common_dist}

\begin{definition}[Bernoulli Distribution]
\label{def:bernoulli}
A random variable $X$ has a \emph{Bernoulli distribution} with parameter $p \in [0,1]$, written $X \sim \text{Bernoulli}(p)$, if $\Prob(X = 1) = p$ and $\Prob(X = 0) = 1 - p$. Then $\E[X] = 1 \cdot p + 0 \cdot (1-p) = p$.
\end{definition}

\begin{definition}[Categorical Distribution]
\label{def:categorical}
A random variable $X$ taking values in $\{1, 2, \ldots, k\}$ has a \emph{categorical distribution} with parameters $(p_1, \ldots, p_k)$, where $p_i \geq 0$ and $\sum_i p_i = 1$, if $\Prob(X = i) = p_i$.
\end{definition}

\begin{definition}[Geometric Distribution]
\label{def:geometric}
$X \sim \text{Geometric}(p)$ for $p \in (0,1]$ if $\Prob(X = k) = (1-p)^{k-1} p$ for $k = 1, 2, \ldots$. Then $\E[X] = 1/p$.
\end{definition}

\begin{remark}
The variance of these distributions ($\Var(X) = p(1-p)$ for Bernoulli, $\Var(X) = (1-p)/p^2$ for Geometric) will be computed in Section~\ref{sec:variance} once variance is formally defined.
\end{remark}

\begin{rlconnection}[Connection to RL]
The categorical distribution models action selection under a stochastic policy: $A_t \sim \pi(\cdot \mid s)$ is categorical with $\Prob(A_t = a) = \pi(a \mid s)$. The geometric distribution arises in episode length analysis: in an MDP with a per-step termination probability $p$, the episode length is geometric.
\end{rlconnection}

\subsection{Joint Distributions and Marginals}
\label{subsec:joint}

\begin{definition}[Joint Distribution]
\label{def:joint}
For random variables $X$ and $Y$, the \emph{joint PMF} (in the discrete case) is:
\begin{equation}
    p_{X,Y}(x, y) = \Prob(X = x, Y = y).
    \label{eq:joint_pmf}
\end{equation}
\end{definition}

\begin{definition}[Marginal Distribution]
\label{def:marginal}
The \emph{marginal PMF} of $X$ is obtained by summing the joint PMF over all values of $Y$:
\begin{equation}
    p_X(x) = \Prob(X = x) = \sum_y \Prob(X = x, Y = y) = \sum_y p_{X,Y}(x, y).
    \label{eq:marginal}
\end{equation}
\end{definition}

\begin{proposition}[Marginalization via Total Probability]
\label{prop:marginalization}
Let $X$ and $Y$ be discrete random variables. Then for every $x$ in the range of $X$:
\[
    p_X(x) = \sum_y p_{X,Y}(x,y).
\]
\end{proposition}

\begin{proof}
Let $Y$ take values in the countable set $\{y_1, y_2, \ldots\}$. The events $B_j = \{Y = y_j\}$ are pairwise disjoint (since $Y$ assigns a single value to each $\omega$), and they cover $\Omega$: every $\omega \in \Omega$ satisfies $Y(\omega) = y_j$ for exactly one $j$, so $\bigcup_j B_j = \Omega$.

Fix an arbitrary $x$ in the range of $X$ and let $A = \{X = x\}$. Since $\bigcup_j B_j = \Omega$:
\[
    A = A \cap \Omega = A \cap \bigcup_j B_j = \bigcup_j (A \cap B_j).
\]
The sets $A \cap B_j = \{X = x, Y = y_j\}$ are pairwise disjoint (inherited from the $B_j$). By countable additivity of $\Prob$ (Definition~\ref{def:prob_measure}):
\begin{align}
    p_X(x) = \Prob(X = x) &= \Prob\!\Bigl(\bigcup_j (A \cap B_j)\Bigr) \nonumber\\
    &= \sum_j \Prob(A \cap B_j) \nonumber\\
    &= \sum_j \Prob(X = x, Y = y_j) = \sum_j p_{X,Y}(x, y_j). \label{eq:marginalization_proof}
\end{align}

Note that this derivation uses only countable additivity and requires no assumption that $\Prob(Y = y_j) > 0$: terms with $\Prob(Y = y_j) = 0$ simply contribute zero to the sum.
\end{proof}

\subsection{Linearity of Expectation}
\label{subsec:linearity}

\begin{keyresult}[Linearity of Expectation]
\begin{theorem}[Linearity]
\label{thm:linearity}
For any discrete random variables $X$ and $Y$ with finite expectations, and constants $a, b \in \R$:
\begin{equation}
    \E[aX + bY] = a\,\E[X] + b\,\E[Y].
    \label{eq:linearity}
\end{equation}
This holds regardless of whether $X$ and $Y$ are independent.
\end{theorem}
\end{keyresult}

\begin{proof}
Let $X$ take values $\{x_i\}$ and $Y$ take values $\{y_j\}$. Denote the joint PMF by $p_{X,Y}(x_i, y_j) = \Prob(X = x_i, Y = y_j)$. Then:
\begin{align}
    \E[aX + bY]
    &= \sum_{i}\sum_{j} (a x_i + b y_j)\, p_{X,Y}(x_i, y_j) \nonumber\\
    &= a \sum_{i}\sum_{j} x_i\, p_{X,Y}(x_i, y_j) + b \sum_{i}\sum_{j} y_j\, p_{X,Y}(x_i, y_j) \nonumber\\
    &= a \sum_{i} x_i \underbrace{\sum_{j} p_{X,Y}(x_i, y_j)}_{= p_X(x_i)} + b \sum_{j} y_j \underbrace{\sum_{i} p_{X,Y}(x_i, y_j)}_{= p_Y(y_j)} \nonumber\\
    &= a \sum_{i} x_i\, p_X(x_i) + b \sum_{j} y_j\, p_Y(y_j) \nonumber\\
    &= a\,\E[X] + b\,\E[Y]. \label{eq:linearity_proof}
\end{align}
The marginal sums $\sum_j p_{X,Y}(x_i, y_j) = p_X(x_i)$ and $\sum_i p_{X,Y}(x_i, y_j) = p_Y(y_j)$ follow from the definition of marginal distributions. No independence assumption is needed.
\end{proof}

\begin{corollary}[Finite Sums]
\label{cor:linearity_sum}
For any random variables $X_1, \ldots, X_n$ with finite expectations and constants $a_1, \ldots, a_n$:
\begin{equation}
    \E\!\left[\sum_{k=1}^n a_k X_k\right] = \sum_{k=1}^n a_k\,\E[X_k].
    \label{eq:linearity_sum}
\end{equation}
\end{corollary}

\begin{proof}
Apply Theorem~\ref{thm:linearity} inductively $n-1$ times.
\end{proof}

\begin{tdconnection}[Connection to TD Learning]
Linearity is used to split the return $G_t = R_{t+1} + \gamma G_{t+1}$ inside expectations:
\[
    \E_\pi[G_t \mid S_t = s] = \E_\pi[R_{t+1} \mid S_t = s] + \gamma\,\E_\pi[G_{t+1} \mid S_t = s].
\]
This is the first step of every Bellman equation derivation.
\end{tdconnection}

%% ============================================================
%% SECTION 4: CONDITIONAL EXPECTATION AND TOWER PROPERTY
%% ============================================================


\end{document}